\addcontentsline{toc}{chapter}{Resumo}

\begin{center}
\huge{{\bf Resumo}}
\vspace{2cm}
\end{center}

% TODO: escrever por ultimo, apos os resultados definitivos.
% Estrutura (max. 2 paragrafos): objetivos + local; metodologia; resultados.

\begin{sloppypar}
[Primeiro parágrafo: contextualizar o problema (automodelo operando em malha
aberta), o objetivo (instrumentar o veiculo para viabilizar controle de
aceleracao em malha fechada e estimacao indireta de torque a partir da corrente)
e o local de realizacao (Laboratorio de Sistemas Inteligentes -- LSI, Escola de
Engenharia da UFMG).]
\end{sloppypar}

\begin{sloppypar}
[Segundo paragrafo: metodologia em uma frase (ensaio de tracao com veiculo
impedido, sensor de corrente ACS758, balanca de mala como referencia de forca)
e resultados principais em numeros: faixa de corrente caracterizada, curva
T = a*I + b obtida, desvio em relacao ao modelo teorico do powertrain.]
\end{sloppypar}

\clearpage
\thispagestyle{empty}
\cleardoublepage
